% =====================================================================
% THE BOUNDARY OF GRAPH SPARSIFICATION FOR COMMUNITY DETECTION
% Version 3 : written clean, 2026-07-26.
%
% RULE FOR THIS FILE (Mohammad, 2026-07-26): no text is inherited from
% main-tnse.tex (v1) or main-v2.tex (v2). Every sentence is written fresh
% against the evidence in PAPER_RESTRUCTURE/. Prior drafts are not
% consulted for prose, framing, or ordering : only the CSVs and the
% verified claims table (phase2_writing/STORY.md) are inputs.
%
% ORDER OF WRITING (Mohammad's rule): introduction first : it carries the
% story. Abstract and conclusion are written LAST.
% =====================================================================

\documentclass[journal]{IEEEtran}

\usepackage{amsmath,amssymb,amsthm}
\usepackage{graphicx}
\usepackage{booktabs}
\usepackage{pgfplots}
\pgfplotsset{compat=1.18}
\usepackage{url}
\usepackage[hidelinks]{hyperref}

\newtheorem{theorem}{Theorem}
\newtheorem{proposition}{Proposition}
\newtheorem{lemma}{Lemma}
\newtheorem{corollary}{Corollary}
\theoremstyle{definition}
\newtheorem{definition}{Definition}
\theoremstyle{remark}
\newtheorem{remark}{Remark}

\begin{document}

\title{The Boundary of Graph Sparsification\\for Community Detection}

\author{Mohammad Dindoost, Bavan DivaaniAazar, Ioannis Koutis, and David A.~Bader%
\thanks{The authors are with the Department of Computer Science,
New Jersey Institute of Technology, Newark, NJ, USA.
E-mail: \{md724, bd344, ikoutis, bader\}@njit.edu}}

\maketitle

%
% ABSTRACT : WRITTEN LAST. Do not draft here until the body is settled.
%
\begin{abstract}
% TODO (last): written after the body. Must state the boundary (degree
% threshold + algorithm family), the cost of the missing control
% (45/63 sign flips), and the protocol : in that order.
\end{abstract}

\begin{IEEEkeywords}
Graph sparsification, community detection, modularity, evaluation methodology,
graph reduction.
\end{IEEEkeywords}

%
\section{Introduction}\label{sec:intro}
% ============================================================
% INTRODUCTION: v3, written clean 2026-07-26.
% Beats B1-B7 per PAPER_RESTRUCTURE/phase4_v3/INTRO_PLAN.md.
% Settled: no contemporary papers named here (D1); no reference to our own
% earlier unpublished draft (D2); dissociation previewed as an open
% phenomenon only (D3); no length pre-compression (D4); mechanism gets one
% line (D5).
% STYLE: no em-dashes, no AI-characteristic phrasing (see STYLE.md).
% ============================================================

% [B1: the practical question]
Community detection is a fundamental task in network
analysis~\cite{fortunato2010community, fortunato202220}. Given a graph, the goal
is to partition its vertices into groups that are densely connected internally
and sparsely connected to the rest of the graph, and the resulting partitions are
used to summarize structure, to identify functional units, and to organize
downstream computation. The task is computationally demanding, and the most
widely used methods, including modularity
optimization~\cite{newman2004finding, blondel2008fast, traag2019louvain} and
flow-based approaches~\cite{rosvall2008maps}, have running times that scale with
the number of edges rather than with the number of vertices. On the graphs to
which these methods are now routinely applied, that quantity is
large~\cite{leskovec2009community}.

Reducing the vertex set is not available as a remedy, because a community
assignment is required for every vertex; removing vertices alters the problem
rather than reducing it. Reducing the edge set is available. Graph
sparsification constructs a subgraph on the same vertex set with substantially
fewer edges while preserving properties of the
original~\cite{spielman2008graph, benczur2015randomized}. If the properties
preserved include community structure, then detection can be performed on the
sparsified graph at lower cost and without loss of quality. That premise has
been pursued for fifteen years, beginning with Satuluri et
al.~\cite{satuluri2011local}, and the resulting methods are now implemented in
standard graph libraries.

% [B2: the answer the field has, and how it was actually scoped]
Satuluri et al.~\cite{satuluri2011local} introduced local similarity
sparsification, in which each vertex retains the incident edges whose endpoints
have the highest Jaccard similarity of neighbourhoods. Evaluating four clustering
algorithms on seven networks, they reported speedups ``often in the range 10-50,
with little or no deterioration in the quality of the resulting clusters,'' and
stated that ``for at least two of the four clustering algorithms, our
sparsification consistently enables higher clustering accuracies.''

The conditions under which that result was obtained are specific. The two
algorithms for which accuracy improved are fixed-$k$, balance-constrained cut
partitioners. Accuracy is measured as agreement with external ground-truth
communities rather than by any objective computed on the graph. The reported
speedups are relative to clustering runs requiring hours on graphs of up to
$1.2 \times 10^8$ edges. The networks on which the gains are largest have average
degrees of 65, 76 and 94, and two of them are derived from high-throughput assays
known to contain false-positive edges. The same paper reports a synthetic sweep
in which the method ``actually outperforms the original clustering starting from
degree 50.'' To our knowledge, no subsequent work examines that threshold.

% [B3: the scope drift, stated without accusation]
Subsequent work extended the recommendation beyond these conditions.
Sparsify-then-detect has since been applied to modularity optimizers as well as
to cut partitioners, to sparse graphs as well as to dense ones, and evaluated
against graph-internal objectives as well as against external references. The
degree threshold reported in the original sweep does not appear in this later
literature, and the conditions attached to the original result are not generally
restated with it.

% [B4: the lost control and what its absence costs (THE HINGE)]
A measurement requirement is also frequently omitted, and the omission has a
direct bearing on what can be concluded. A sparsifier intended for community
detection removes edges lying between communities in preference to edges lying
within them. Quality measures for a partition, including modularity and
conductance, penalize exactly the between-community edges that such a sparsifier
removes. A partition evaluated on the sparsified graph is therefore evaluated
against a graph from which its penalty term has been partially deleted, and the
resulting quantity reflects the sparsifier's selectivity rather than the extent
to which community structure was preserved. The comparison that answers the
intended question evaluates the partition obtained from the sparsified graph and
the partition obtained from the original graph on the same graph, namely the
original one.

This requirement is stated in the 2011 paper, whose authors write that they
``cannot simply measure the conductance using the very same sparsified graph,
since that would tell us nothing about how well the sparsified graph retained the
cluster structure in the original graph,'' and who report every quality figure on
the original graph accordingly. It is absent from much of the work that followed,
and from the released evaluation code of a widely used sparsification benchmark.

The effect of the omission is large. Across 63 controlled configurations spanning
four detection algorithms, two sparsifier families and seven networks, evaluating
on the sparsified graph rather than on the original reverses the sign of the
conclusion in 45 of them. In the most extreme configuration, a pipeline that
reduces modularity by 0.40 under the correct evaluation appears to increase it by
0.39 under the incorrect one. Under such an evaluation a boundary cannot be
located, since improvement is reported throughout the parameter range.

% [B5: what we did]
We therefore rebuilt the evaluation before re-asking the question. Partitions are
always scored on the original graph. Comparisons against a stochastic baseline
are given the same wall-clock budget as the pipeline they are compared with, and
are also checked against a fixed-restart baseline, because a runtime-matched
control is weak precisely when the pipeline is fast. Because sparsification
fragments partitions, and because best-match and information-theoretic recovery
scores rise mechanically with the number of clusters, every recovery comparison
is matched for granularity against a partition of the original graph at the same
cluster count, and is referenced to a chance floor computed from size-matched
random partitions. Because a heavy-tailed degree sequence alone can produce
several of the effects attributed to community structure, the central mechanistic
claims are checked against degree-preserving randomizations of the same graphs.

With that instrument we ran 23 controlled experiments over both sparsifier
families that have ever claimed quality gains, degree-based and similarity-based,
the latter being the 2011 method itself; four detection algorithms drawn from two
families; seventeen real networks of up to 25 million edges; and a synthetic
sweep of the one variable the original authors identified as decisive. The design
locates a boundary rather than confirming or refuting a claim, and it is reported
as such. Predictions and stopping rules were fixed before each experiment ran,
and several of our own hypotheses did not survive them.

% [B6: the boundary]
Our results locate a boundary, and the transition across it is abrupt.

On one side of it sparsification does not improve detection. For modularity-family
detectors with free granularity, on graphs with average degree between 5 and 33,
no sparsifier we tested improves modularity measured on the original graph beyond
the variation of a small number of random restarts, at any retention level that
preserves quality, under any of the four algorithms examined. Nor is there a
speed benefit. Measured end to end, including the cost of the sparsifier itself,
the pipeline requires between $0.87\times$ and $1.35\times$ the time of clustering
the original graph directly, and at aggressive retention it is slower, because a
fragmented graph presents the optimizer with more work rather than less. Every apparent
gain we found in this regime, in modularity, in ground-truth recovery, and in the
one case where a sparsifier's edge-selection signal is demonstrably
structure-aware rather than degree-driven, dissolves when granularity is matched,
when chance is subtracted, or when the baseline is given an equal compute budget.

On the other side of the boundary sparsification does improve detection, under
the conditions the original authors described. On synthetic graphs with planted communities, holding the number of
communities fixed by construction so that no granularity effect is possible,
similarity-based sparsification raises both the objective measured on the
original graph and chance-corrected agreement with the planted labels once the
average degree passes approximately 50. The effect is absent at average degree 11
and 25, present in two thirds of configurations at 49, and present in all of them
at 102 and 221. It survives the worst case over ten independent runs of the
partitioner, and it is not reproduced by degree-based or uniform sparsification
at identical retention, so it is a property of the similarity signal rather than
of the edge budget. Injecting spurious edges makes the gain grow monotonically,
which is the signature the denoising explanation predicts.

What decides the outcome is not which sparsifier is used. It is the average
degree of the graph, the family of the detection algorithm, since fixed-$k$
balanced partitioners gain where free-granularity modularity optimizers do not,
and whether quality is measured against an external reference or against an
objective defined on the graph whose edges were deleted. The mechanism is
consistent with this. The quantity governing a degree-based sparsifier's effect
on modularity is the separation between the scores it assigns to within-community
and between-community edges, and we show by direct intervention that this
quantity is controlled by where high-degree-product edges sit relative to
community boundaries, at fixed degree sequence, fixed partition and fixed
modularity.

Two qualifications attach to this result. Sparsification repairs a weaker
detector rather than surpassing a stronger one: a resolution-tuned modularity
optimizer on the untouched graph outperforms every sparsified pipeline we
measured. And with an exact similarity computation the quality gain is not a
speed gain, because the sparsifier costs more than the detection it accelerates.

% [B7: contributions]
This paper contributes the following.

\begin{itemize}
  \item \textbf{The boundary.} A degree threshold near average degree 50 and an
  algorithm-family split, established under controls that exclude granularity by
  construction, together with the negative territory delimited across two
  sparsifier families, four detection algorithms and seventeen networks.

  \item \textbf{An evaluation protocol, and the measured cost of omitting it.}
  Each control is motivated by a specific failure it catches, and the consequence
  of the central omission is quantified at 45 sign reversals in 63
  configurations. We also introduce a permuted-weight control for
  similarity-weighted pipelines, which shows that the benefit attributed to
  similarity weighting is reproduced, and in one case exceeded, when the weights
  are randomly permuted across edges.

  \item \textbf{A mechanism for the boundary.} A decomposition of the
  score-separation quantity into a sorting term and a granularity term across
  seventeen networks, and a direct causal test in which that quantity is steered
  through more than a full sign change while degree sequence, partition,
  intra-edge fraction and modularity are all held exactly fixed.

  \item \textbf{An open phenomenon.} In several controlled settings the
  partitions that best recover reference communities are not the partitions that
  score best on the objective, and the two criteria move in opposite directions
  within the same configuration. We report this as an observation rather than a
  result, and argue that it is the most consequential open question the boundary
  exposes.
\end{itemize}


%
% ORDER SETTLED 2026-07-26, after the introduction. It follows B6: the
% negative territory, then the boundary, then the mechanism that explains
% both. This replaces the provisional skeleton, which had the mechanism
% before the negative territory. Mohammad: the order may still change, and
% sections may be merged or moved once each is written.
%

\section{Background and Related Work}\label{sec:background}
% ============================================================
% SECTION 2: BACKGROUND AND RELATED WORK. Written 2026-07-27.
%
% Job, per PAPER_STRUCTURE.md and the HANDOFF: introduce only what the rest of
% the paper uses; define "removable redundancy", which Section 1 uses as a
% technical term and which nothing else in the paper defines; and hold the
% contemporary instances of the missing control, with the arithmetic (decision
% D1). The literature is motivation and related work here, never the spine.
%
% Mohammad, 2026-07-27: Socievole and Pizzuti are cited as current practice,
% briefly, with no tables and no numbers. This is a related-work section.
%
% Every characterization below comes from a paper read in full, and the reading
% is recorded in PAPER_RESTRUCTURE/references/:
%   NOTES_satuluri2011.md          the 4.3 transfer control and the fixed-k arms
%   NOTES_chen_and_fastcd.md       Chen (what they own, what they do not) and
%                                  Laeuchli (the regime where sparsification pays)
%   NOTES_effres_2026.md           Pari et al., football 0.87 against max 0.6046
%   NOTES_kcore_ijain2026.md       Setiadi et al., karate 0.63 against max 0.4198
%   NOTES_socievole_pizzuti_2026.md  read in full; cited here as practice only
%   P1_EVIDENCE.md                 Hamann, Blagus, sampling line, IJCAI survey
%   P2_BACKGROUND_RESEARCH.md      metric backbone, Maslov-Sneppen, Graclus/MLR-MCL
% STYLE: no em-dashes, no AI-characteristic phrasing (see STYLE.md).
% ============================================================

\subsection{Detection algorithms and their degrees of freedom}

Community detection algorithms divide into two families, and the division is the
first of the two conditions this paper identifies.

One family chooses its own number of communities. Modularity
optimizers~\cite{newman2004finding} such as Louvain~\cite{blondel2008fast} and
Leiden~\cite{traag2019louvain} maximize \eqref{eq:modularity} by agglomeration,
stopping when no move improves the objective. Infomap~\cite{rosvall2008maps}
minimizes the description length of a random walk, and label
propagation~\cite{raghavan2007near} iterates a local majority rule to a fixed
point. None takes the number of communities as an input, and each returns
whatever count its objective and its search prefer. Modularity in particular has
a resolution limit~\cite{fortunato2007resolution}: at $\gamma = 1$ it cannot
resolve communities below a size set by the total number of edges, so the count
it returns is a property of the graph's size as much as of its structure.

The other family takes the number of communities as an input. Multilevel graph
partitioners such as Metis~\cite{karypis1998metis} and
Graclus~\cite{dhillon2007weighted} coarsen the graph, partition the coarsest
level into a requested $k$ parts, and refine on the way back up, subject to a
balance constraint on part sizes. Flow-based clustering with a coarsening
schedule, as in MLR-MCL~\cite{satuluri2009scalable}, sits between the two: the
count follows from an inflation parameter rather than from a requested $k$.

The difference matters here because sparsification changes what a
free-granularity algorithm returns, and does not change what a fixed-$k$
algorithm returns. Section~\ref{sec:protocol} makes that a control and
Section~\ref{sec:experiments} makes it a result.

\subsection{Sparsification and what it is guaranteed to preserve}

Spectral sparsification is the part of this literature with proofs. Benczúr and
Karger~\cite{benczur2015randomized} preserve all cuts to within $1 \pm \epsilon$
by sampling edges according to connectivity, and Spielman and
Srivastava~\cite{spielman2008graph} preserve the Laplacian quadratic form by
sampling according to effective resistance. The guarantee is about the Laplacian,
not about communities, and Section~\ref{sec:intro} states why the two are related
without being the same.

The methods actually used before community detection are cheaper and carry no
such guarantee. Similarity-based methods rank each edge by the overlap of its
endpoints' neighbourhoods, as in the local sparsification of Satuluri et
al.~\cite{satuluri2011local}. Degree-based methods score an edge from its
endpoints' degrees alone, as in DSpar~\cite{liu2023dspar}, which approximates
effective resistance by $1/d_u + 1/d_v$. Backbone extraction keeps edges that are
statistically surprising against a local null, as in the disparity
filter~\cite{serrano2009extracting} and its noise-corrected
successor~\cite{coscia2017backboning}. Learned methods train a model to choose the
subgraph~\cite{wu2020gsgan}. Hamann et al.~\cite{hamann2016structure} compared
these families systematically on social networks and released the
implementations that are now distributed in standard graph libraries.

Two recent surveys bound the field. Chen et al.~\cite{chen2024demystifying}
benchmark twelve sparsifiers against sixteen graph properties on fourteen
networks and conclude that no single sparsifier preserves all of them and that
the choice should follow the downstream task. They also establish, at benchmark
scale, that pruning fragments graphs and that fidelity to the full graph's
partition degrades monotonically with the prune rate. Both facts are theirs and
this paper does not claim them. Hashemi et al.~\cite{hashemi2024survey} survey
the wider reduction field of sparsification, coarsening and condensation, and
their task list is a useful negative: the modern literature has reoriented around
graph neural network prediction, and community detection has largely dropped out
of it as an evaluation target.

\subsection{Removable redundancy}\label{sec:background:redundancy}

The second condition in this paper's claim is that a graph carry edges whose
removal does not destroy its community structure. The quantity has an established
formal relative. The \emph{metric backbone} of a weighted graph is the union of
its all-pairs shortest paths, so every edge outside it is redundant in the exact
sense that no shortest path uses it. Correia et
al.~\cite{correia2023contact} report that the backbone is a small fraction of the
edges of real contact networks, and Dreveton et al.~\cite{dreveton2024metric}
prove that deleting everything outside it leaves community structure intact on a
broad class of weighted random graphs with communities.

Removable redundancy in the sense used here is broader than the metric one, since
the edges a sparsifier may delete without cost are not only those absent from
shortest paths, and we do not have a procedure that measures it directly. What
the results of Section~\ref{sec:experiments:boundary} establish is that the
quantity varies across graphs, that average degree is its strongest measured
correlate, and that degree heterogeneity is not a correlate at all.

\subsection{Sparsify, then detect}

Satuluri et al.~\cite{satuluri2011local} introduced the pipeline this paper
examines and reported speedups commonly between $10\times$ and $50\times$ with
little or no loss of cluster quality, together with improved accuracy for two of
the four clustering algorithms they tested. Those two were the partitioners that
take $k$ as an input.

The practice spread. Reference implementations entered graph
libraries~\cite{hamann2016structure}, sparsification was applied before community
detection at industrial scale~\cite{satuluri2020simclusters}, generative models
were trained to produce the subgraph~\cite{wu2020gsgan}, and triangle-aware
spectral sparsifiers were proposed specifically because sparsification appeared
to be a route to community detection~\cite{sotiropoulos2021triangle}. It remains
current: Socievole and Pizzuti weight a graph by effective resistance, apply a
kernel, sparsify, and then detect communities on what
remains~\cite{socievole2024community, socievole2026effective}, and further
effective-resistance and $k$-core variants of the same pipeline have appeared
since~\cite{pari2026effective, setiadi2025community}.

A parallel line asks whether sampling preserves community structure at all.
Leskovec and Faloutsos~\cite{leskovec2006sampling} set the sampling problem,
Maiya and Berger-Wolf~\cite{maiya2010sampling} sample explicitly for community
structure, and Blagus et al.~\cite{blagus2015sampling} report that sampling can
manufacture the appearance of community structure rather than preserve it. Peng
et al.~\cite{peng2014accelerating} reduce the graph by $k$-core before detecting,
which is the coarsening-flavoured version of the same idea and removes vertices
as well as edges.

\subsection{How the pipeline is evaluated}

The evaluation question that Section~\ref{sec:protocol} formalizes was raised in
the paper that introduced the pipeline. Satuluri et al.\ state in their \S4.3
that the clusters obtained from a sparsified graph cannot be scored on that same
graph, because doing so says nothing about how well the sparsification retained
the structure of the original, and they report every conductance on the original
graph accordingly. They also fix the number of communities once per network and
supply it to every arm, so their comparisons are matched on granularity by
construction, and they charge the sparsification time to the speedup. On all
three counts the original paper is stricter than much of what followed it.

What followed is less uniform. In the released code of the most thorough public
benchmark of sparsification, modularity for the sparsified condition is computed
on the sparsified graph~\cite{chen2024demystifying}. Two recent
sparsify-then-detect papers report objective values that exceed what their graphs
admit: an effective-resistance method reports modularity $0.87$ on the American
College Football network~\cite{girvan2002community}, whose maximum over all
partitions is $0.6046$, and a $k$-core method reports $0.63$ on Zachary's karate
club~\cite{zachary1977information}, whose maximum is
$0.4198$~\cite{pari2026effective, setiadi2025community}. A modularity above the
graph's maximum is not a matter of interpretation; it identifies the graph being
scored as the sparsified one. Section~\ref{sec:experiments:accounting} measures
what that choice is worth.

Two further hazards are documented prior art rather than findings of this paper.
Hamann et al.~\cite{hamann2016structure} warn that normalized mutual information
reports higher similarity for larger numbers of communities, which makes it
unsuitable for comparing partitions of sparsified networks, and point to
chance-corrected indices~\cite{vinh2010information} in its place. Gottesb\"{u}ren
et al.~\cite{gottesburen2025linear} observe that the cost of sparsifying can
cancel the speedup it is meant to produce. Separating a degree-driven effect from
a community-driven one requires a null that holds the degree sequence fixed, for
which the standard instrument is the degree-preserving edge
swap~\cite{maslov2002specificity} over configuration-model
graphs~\cite{molloy1995critical}.

Finally, one paper marks the boundary of this paper's scope rather than falling
inside it. Laeuchli~\cite{laeuchli2020fast} analyses community detection on
sparsified graphs for spectral solvers on dense planted-partition models, which
is the regime in which sparsification genuinely pays, because the detection cost
there is superlinear in the number of edges and the planted structure survives
heavy deletion. The graphs and detectors examined here are the ones the pipeline
is applied to in practice, and they are not that regime.


\section{Evaluation Protocol}\label{sec:protocol}
% ============================================================
% SECTION 3: EVALUATION PROTOCOL. Written clean 2026-07-26.
% Title chosen by Mohammad 2026-07-26 over "What Honest Evaluation Requires":
%   "honest" implies the alternative is dishonesty, which is not this paper's
%   position. Satuluri et al. stated the requirement correctly in 2011.
% Weight: MEDIUM-HEAVY (PAPER_STRUCTURE.md). This is the instrument.
%
% Evidence, control by control:
%   transfer scoring   TERMINOLOGY.md; exp_P SUMMARY V2 (45/63, 60/63, 15/63);
%                      exp_G SUMMARY Table A (uniform, 6/6 sign flips)
%   granularity        exp_P finding 4 (com-DBLP Louvain 0.1930 -> 0.2996 at
%                      gamma=576); exp_L (Leiden 0.1930 / 0.2971)
%   chance floor       exp_P finding 4 (com-Amazon 0.0086@k=231 -> 0.0567@k=8651)
%   compute matching   exp_P V1/V3 and caveat C3; exp_C section 1 (2 restarts
%                      in all 48 cells); exp_AA (budget buys 1 restart)
%   configuration null exp_B SUMMARY (17 networks; dQ_fixed and delta tables)
%   permuted weights   exp_V V2 (email-Enron; ca-CondMat rewiring)
%   realized retention exp_C section 0 (sampler audit); exp_AB section 1 (floors)
%   end-to-end cost    exp_P V5 (0.66x max among quality-preserving cells)
%
% NOTE: the modularity definition sits here because the m-dependence of its
% null term IS the argument in 3.1. If Section 2 later defines it, delete
% Eq. (1) here and cross-reference instead.
% STYLE: no em-dashes, no AI-characteristic phrasing (see STYLE.md).
% ============================================================

Section~\ref{sec:intro} listed four requirements, each belonging to a specific
class of method rather than to sparsification in general. This section states
them as controls. Each is given with the failure it catches and with what we
measure when it is absent. Table~\ref{tab:protocol} collects them. Every number
in the sections that follow is produced under all of them.

\subsection{One yardstick for two partitions}\label{sec:protocol:transfer}

Detection may use any graph; scoring must use the original one. Let
$\mathcal{A}$ be a detection algorithm, $G$ the original graph and $G'$ a
sparsified copy of it. The two partitions under comparison are
$P_{\mathrm{sparse}} = \mathcal{A}(G')$ and $P_{\mathrm{base}} = \mathcal{A}(G)$,
and since both are partitions of the same vertex set, the question the pipeline
raises is which of them better describes $G$. Two accountings of that question
appear in this literature:
\begin{align}
\Delta Q_{\mathrm{transfer}} &= Q(P_{\mathrm{sparse}}, G) - Q(P_{\mathrm{base}},
G), \label{eq:transfer}\\
\Delta Q_{\mathrm{self}} &= Q(P_{\mathrm{sparse}}, G') - Q(P_{\mathrm{base}}, G).
\label{eq:self}
\end{align}
They differ in one place, the graph on which the first term is evaluated.
Equation~\eqref{eq:transfer} measures both partitions with a single instrument.
Equation~\eqref{eq:self} changes the instrument between the arm under test and
the arm it is tested against, and the instrument it applies to the former is the
graph from which edges have been removed. We report
$\Delta Q_{\mathrm{transfer}}$ throughout; $\Delta Q_{\mathrm{self}}$ appears in
released benchmark tooling and in recent work in this line, documented in
Section~\ref{sec:related}.

The two accountings diverge systematically, and the objective shows why.
Modularity scores a partition $P$ of a graph with $m$ edges by
\begin{equation}\label{eq:modularity}
Q(P, G) \;=\; \sum_{C \in P} \left[ \frac{m_C}{m} - \gamma
\left( \frac{D_C}{2m} \right)^{\!2} \right],
\end{equation}
with $m_C$ the number of edges inside community $C$, $D_C$ the sum of degrees in
$C$, and $\gamma$ a resolution parameter equal to one in the standard definition.
The subtracted term is normalized by $m$. Deleting edges shrinks it, so
modularity computed on a sparsified graph rises whether or not the partition has
improved, and the rise is larger the more edges are removed.

The effect is not small and it is not confined to sparsifiers that select
intelligently. Table~\ref{tab:uniform} scores uniform random edge deletion, whose
selection rule uses no information at all, both ways on three networks at two
retention levels. The sign of the conclusion is opposite in every one of the six
cells, and the gap between the two accountings roughly doubles when retention
falls from $0.9$ to $0.8$, which is the behaviour \eqref{eq:modularity} predicts.

\begin{table}[t]
\centering
\caption{Uniform random edge deletion under the two accountings, Leiden, mean of
three seeds. $\Delta Q_{\mathrm{self}}$ is positive in every cell and
$\Delta Q_{\mathrm{transfer}}$ is negative in every cell. Uniform deletion uses
no structural information, so the inflation comes from the objective rather than
from any selection rule.}
\label{tab:uniform}
\small
\setlength{\tabcolsep}{4pt}
\begin{tabular}{@{}lccrr@{}}
\toprule
Network & $Q_{\mathrm{base}}$ & $\rho$ & $\Delta Q_{\mathrm{self}}$ &
$\Delta Q_{\mathrm{transfer}}$\\
\midrule
ca-CondMat  & 0.729 & 0.800 & $+0.0111$ & $-0.0164$\\
ca-CondMat  & 0.729 & 0.900 & $+0.0053$ & $-0.0064$\\
email-Enron & 0.611 & 0.800 & $+0.0052$ & $-0.0162$\\
email-Enron & 0.611 & 0.900 & $+0.0011$ & $-0.0086$\\
com-DBLP    & 0.826 & 0.800 & $+0.0078$ & $-0.0247$\\
com-DBLP    & 0.826 & 0.900 & $+0.0040$ & $-0.0098$\\
\bottomrule
\end{tabular}
\end{table}

Across a larger grid the pattern holds under every algorithm we tested.
Table~\ref{tab:accounting} covers 63 configurations spanning three detection
algorithms, seven networks and three sparsifier settings. Modularity read off the
sparsified graph is positive in 60 of them and reaches $+0.537$; the same
partitions scored on the original graph are positive in 15. The sign of the
conclusion differs between the two accountings in 45 of the 63 configurations.
Eleven cells exceed a compute-matched baseline, all of them under label
propagation, and they are not a subset of the 15: two of them sit below the plain
baseline mean while beating a matched control that had drawn too few restarts,
which is the failure Section~\ref{sec:protocol:compute} treats.

\begin{table}[t]
\centering
\caption{The two accountings across 63 configurations, seven networks by three
sparsifier settings under each algorithm. A flip is a cell where
$\Delta Q_{\mathrm{self}} > 0$ and $\Delta Q_{\mathrm{transfer}} < 0$, so that
the two accountings disagree on the sign of the conclusion.}
\label{tab:accounting}
\small
\setlength{\tabcolsep}{4pt}
\begin{tabular}{@{}lcccc@{}}
\toprule
 & Infomap & Louvain & Lab.\ prop. & Total\\
\midrule
Cells                              & 21 & 21 & 21 & 63\\
$\Delta Q_{\mathrm{self}} > 0$     & 20 & 20 & 20 & 60\\
$\Delta Q_{\mathrm{transfer}} > 0$ & 6  & 0  & 9  & 15\\
Flips                              & 14 & 20 & 11 & 45\\
max $\Delta Q_{\mathrm{self}}$ & $+0.312$ & $+0.260$ & $+0.537$ & $+0.537$\\
\bottomrule
\end{tabular}
\end{table}

\subsection{Matching granularity}\label{sec:protocol:granularity}

Sparsification fragments a graph, and an algorithm that chooses its own
granularity responds by returning more communities. Chen et
al.~\cite{chen2024demystifying} report community counts rising by three orders of
magnitude across a pruning sweep. Purity, normalized mutual information and the
average best-match $F_1$ over reference communities that we report all increase
with the number of communities, so a recovery comparison between partitions of
different granularity measures that difference along with whatever else it is
meant to measure. Hamann et al.~\cite{hamann2016structure} warned that normalized
mutual information is unsuitable for comparing partitions of sparsified networks,
because it reports higher similarity for larger numbers of
communities~\cite{vinh2010information}, and pointed to the adjusted Rand index
instead; the control below operationalizes that warning.

The control is a partition of the \emph{original} graph at a matched community
count, obtained by raising the resolution parameter $\gamma$ in
\eqref{eq:modularity} until the counts agree. Table~\ref{tab:granularity} applies
it to com-DBLP under Louvain. Local similarity sparsification at realized
retention $0.503$ raises average best-match $F_1$ from $0.102$ to $0.193$, an
apparent gain of $+0.091$, while raising the number of communities of size at
least three from 220 to 10,947. The original graph partitioned at $\gamma = 640$
reaches 11,159 such communities and scores $0.300$. The finer partition is worth
more than twice what the sparsification appeared to be worth, and it requires no
sparsifier. The same holds for degree-based sparsification on the same network,
where the apparent change is negative and the matched control is positive.
Leiden reproduces the sparsified figure to four decimal places and its own
control, at $\gamma = 576$ and 10,812 communities, reaches $0.297$, so the effect
belongs to the modularity objective's default granularity rather than to one
optimizer's search.

\begin{table}[t]
\centering
\caption{Granularity control on com-DBLP under Louvain, average best-match $F_1$
over reference communities of size at least three. $k_{\geq 3}$ counts
communities of that size; $\gamma$ is the resolution used on the \emph{original}
graph. Each sparsified arm is followed by its control, a partition of the
original graph at the same $k_{\geq 3}$. The floor is a size-matched random
partition at the arm's own $k_{\geq 3}$. Baseline and sparsified rows are means
over three to six seeds; controls are single runs.}
\label{tab:granularity}
\small
\setlength{\tabcolsep}{4pt}
\begin{tabular}{@{}lrrcc@{}}
\toprule
Condition & $k_{\geq 3}$ & $\gamma$ & $F_1$ & floor\\
\midrule
Original graph               & 220    & 1   & $0.102 \pm 0.005$ & 0.019\\
DSpar, $\rho = 0.500$        & 1{,}684  & 1   & $0.098 \pm 0.002$ & 0.049\\
\quad matched control        & 1{,}657  & 24  & $0.125$           & \\
L-Spar, $\rho = 0.503$       & 10{,}947 & 1   & $0.193 \pm 0.000$ & 0.090\\
\quad matched control        & 11{,}159 & 640 & $0.300$           & \\
\bottomrule
\end{tabular}
\end{table}

Where the counts cannot be brought together, we draw no recovery conclusion. The
comparability rule is $|\Delta k|/k < 25\%$, and the cells it excludes are
reported as excluded rather than compared.

\subsection{Subtracting chance}\label{sec:protocol:chance}

Recovery is reported with chance-corrected indices. Where an uncorrected measure
is used because the reference communities are overlapping and numerous, as with
average best-match $F_1$ on the largest networks, we also compute the score of a
random partition with the same community size distribution, and report it beside
the measurement.

That floor moves with the quantity the previous control governs. On com-Amazon
under Louvain, a size-matched random partition scores $0.009$ against the
reference communities at 231 communities and $0.057$ at 8,651, a $6.6\times$ rise
produced by fragmentation alone. Any comparison of best-match $F_1$ across
different community counts reads part of that slope, and subtracting the floor is
what separates the two.

\subsection{An equal compute budget}\label{sec:protocol:compute}

Every detection algorithm we use is stochastic, so a single run is a draw rather
than an answer, and a pipeline that sparsifies and then detects has spent more
time than one detection run. The registered control gives the baseline as many
independent restarts as fit the pipeline's wall clock and takes the best of them.

That control is weaker than it appears, and its weakness favours the treatment.
When the pipeline is fast the budget buys very little: two restarts in all 48
cells of one sweep, and between 1 and 16 in another, with 1 in most cells on the
two largest networks. Label propagation shows what this costs. Under the
compute-matched control it is the strongest of the three algorithms in that
sweep, with a mean change of $+0.009$; against the best of 20 plain restarts of
the same baseline it is the weakest, at $-0.106$. On email-Enron its two apparent
gains of $+0.146$ and $+0.161$ become $-0.068$ and $-0.053$, because the
compute-matched baseline had drawn 1 to 2 restarts and never sampled the
algorithm's good mode, whose modularity is $0.523$ against the matched control's
$0.308$.

We therefore report both: the compute-matched figure, which is the fair
accounting of what the pipeline costs, and the best of a fixed number of
restarts, which is a strictly more generous baseline that can only handicap the
treatment. A result counts as surviving only if it survives both. This is also
the reason the granularity control above is applied to results that have already
passed the compute test, rather than in place of it.

\subsection{Nulls}\label{sec:protocol:nulls}

Two of the quantities this literature reads as evidence of community structure
can be produced without any community structure at all. Each has a null that
detects the substitution.

\paragraph{A degree sequence without communities}
A degree-based sparsifier selects on a function of the degree sequence, so an
effect it produces on a heavy-tailed graph may follow from that sequence alone.
The control is a degree-preserving rewiring~\cite{molloy1995critical}, built by
double edge swaps, which holds every vertex degree exactly and destroys the
community structure. A signal that reappears on the rewiring is not evidence
about communities.

Two quantities fail this test. The first is the change in modularity of a fixed
partition when the graph is sparsified, which is positive on all seventeen
networks we tested; on their rewirings it is positive on all seventeen as well,
and larger than the real value on fifteen of them. The second is the separation
between the scores a degree-based sparsifier assigns to within-community and
between-community edges, which on the rewirings is larger than on the real graphs
in sixteen of seventeen cases, with a median ratio of $1.60$. Real community
structure suppresses that separation rather than producing it. Neither quantity
is used as evidence in this paper. The mechanism they belong to is the subject of
Section~\ref{sec:mechanism}, where the same null is what makes the causal test
interpretable.

\paragraph{Weights without their placement}
A related family of proposals keeps every edge and reweights it by a similarity
score. The corresponding null permutes the computed weights uniformly at random
across the edges, which preserves the weight distribution exactly and destroys
every association between a weight and the structure that produced it. We are not
aware of its use in this literature.

It changes the reading of the one modularity gain we obtained from four ways of
deploying a similarity signal. On email-Enron, weighting each edge by
$1 + J(u,v)$, with $J$ the Jaccard similarity of the endpoint neighbourhoods,
raises modularity measured on the original graph to $0.6153 \pm 0.0034$ against a
baseline of $0.6071$ and a compute-matched best of $0.6051$, a gain of $+0.0102$
at five pooled standard deviations. Permuting those same weights across edges
gives $0.6168 \pm 0.0020$, a gain of $+0.0117$, which is larger. On a
degree-preserving rewiring of ca-CondMat, where the similarity signal is absent
by construction, weighting still produces $+0.0023$ and clears the same
significance test. What the weights contribute is heterogeneity in the
optimizer's search rather than information about communities.

\subsection{Realized retention}\label{sec:protocol:retention}

Retention is measured, never assumed. The sampling scheme in common use clips
per-edge probabilities at one, and between 13\% and 35\% of edges reach that
clip, so the expected number of retained edges falls below the nominal $\alpha m$
and saturates: at $\alpha = 1.0$ one of our networks retains $0.662$ of its
edges, and at $\alpha = 0.95$ another retains $0.447$. Three samplers that all
describe themselves by the same $\alpha$ realize $0.33$ to $0.55$, $0.37$ to
$0.68$, and, once the scale factor is solved by bisection, $0.70$ to $0.95$.
Sparsifiers are compared at matched realized retention throughout, and the
realized value appears in every table.

Two properties of the local selection rules limit what can be matched. The
achievable retentions form a coarse grid, so that on com-Amazon local similarity
sparsification can produce only $0.301$ and $0.542$ and nothing between them. And
several methods cannot reach aggressive retention on a sparse graph at all,
because each vertex keeps at least one incident edge; the resulting floors reach
$0.36$ on the sparsest network we test, which is reported with the results rather
than treated as a free parameter.

A comparison must also be like for like in what is removed. A method that deletes
vertices as well as edges scores its partition on a smaller and denser
subpopulation, so the fraction of vertices retained belongs in the table beside
the fraction of edges. Every sparsifier in this study removes edges only, and we
report both fractions.

\subsection{Cost measured end to end}\label{sec:protocol:cost}

The sparsifier is a computation and is charged to the pipeline that uses it, as
it was in the original report~\cite{satuluri2011local} and as Gottesb\"{u}ren et
al.~\cite{gottesburen2025linear} recommend after measuring the same overhead in
graph partitioning. The comparison is the whole pipeline, sparsifier included,
against one run of the same detection algorithm on the original graph.

Among the cells of our largest sweep that preserve quality, the greatest
end-to-end speedup is $0.66\times$, which is slower than not sparsifying.
Speedups above $1.5\times$ occur in 12 of 63 cells, and every one of them loses
modularity. Discounting the sparsifier entirely and timing only the detection
step, the median speedup is between $1.15\times$ and $1.63\times$ by algorithm,
well short of the factor of two that halving the edges suggests, because all
three algorithms return more communities on a sparsified graph and do more work
per pass.

\begin{table*}[t]
\centering
\caption{The protocol. Each control is motivated by a specific failure, and the
last column reports what we measure when the control is absent.}
\label{tab:protocol}
\begin{tabular}{@{}p{0.19\textwidth}p{0.29\textwidth}p{0.42\textwidth}@{}}
\toprule
Failure & Control & Measured consequence of omitting it\\
\midrule
Quality read off the sparsified graph &
Score both partitions on the original graph, Eq.~\eqref{eq:transfer} &
The sign of the conclusion differs in 45 of 63 configurations; positive in 60 of
63 one way and in 15 of 63 the other (\S\ref{sec:protocol:transfer})\\
Finer partitions score higher &
Partition the original graph at a matched community count &
com-DBLP: an apparent $+0.091$ in best-match $F_1$ against a control that reaches
$+0.198$ with no sparsifier (\S\ref{sec:protocol:granularity})\\
Uncorrected agreement measures &
Chance-corrected indices, and a size-matched random-partition floor &
The floor rises $6.6\times$ between 231 and 8,651 communities on one graph
(\S\ref{sec:protocol:chance})\\
Unequal compute &
Compute-matched restarts \emph{and} the best of a fixed number &
One algorithm ranks first under the former and last under the latter; two gains
near $+0.15$ become losses (\S\ref{sec:protocol:compute})\\
Degree mechanics read as structure &
Degree-preserving rewiring &
Both quantities are reproduced by the null, and are larger on it in 15 of 17 and
16 of 17 networks (\S\ref{sec:protocol:nulls})\\
Weighting credited to similarity &
Permute the weights across edges &
The permuted weights reproduce the gain and exceed it, $+0.0117$ against
$+0.0102$ (\S\ref{sec:protocol:nulls})\\
Nominal retention &
Report and match on measured retention &
The same nominal $\alpha$ realizes retentions from $0.33$ to $0.95$ depending on
the sampler (\S\ref{sec:protocol:retention})\\
Sparsification charged to nobody &
Time the pipeline, sparsifier included &
No quality-preserving configuration exceeds $0.66\times$
(\S\ref{sec:protocol:cost})\\
\bottomrule
\end{tabular}
\end{table*}

\subsection{Pre-registration}\label{sec:protocol:prereg}

Each experiment in this study has a design document written before its code,
fixing what is measured, what is predicted, and what outcome would end the line of
work. Several of the predictions were ours and failed. The synthetic sweep that
locates the boundary was registered with the prediction that sparsification would
not produce an honest modularity gain at any density, and the fixed-$k$ arm
falsified it.


\section{Experiments}\label{sec:experiments}
% UNDER RECONSTRUCTION (2026-07-26 late): the new experiments section absorbs the
% measurement half of the old protocol section and all of the old "below the
% boundary" section. Both originals are preserved under sections/_mined/.
% ============================================================
% SECTION 4 OPENING + 4.1 SETUP. Written 2026-07-26 (late).
%
% Part of the restructure to Mohammad's shape: introduction, background,
% protocol (argument only), experiments (everything measured), discussion,
% conclusion. This file is the opening and the setup; the result subsections
% follow in 04_below.tex until that is mined and replaced.
%
% Numbers here come from exp_AB/results.csv (the seven networks),
% exp_B/results.csv (the seventeen), exp_AA (LFR), exp_AC (SBM/dcSBM),
% exp_AD (Metis on real graphs). All verified against CSVs, not summaries.
% STYLE: no em-dashes, no AI-characteristic phrasing (see STYLE.md).
% ============================================================

Everything below runs under the protocol of Section~\ref{sec:protocol}. The
experiments answer three questions in order. What does the accounting cost, that
is, how much of the reported benefit in this literature is a property of how it
is measured. Where the pipeline is used in practice, does it help. And where it
does help, what are the conditions.

\subsection{Setup}\label{sec:experiments:setup}

\paragraph{Networks}
The quality measurements use the seven real networks of Table~\ref{tab:networks},
chosen to span an order of magnitude in average degree while remaining tractable
under a protocol that reruns every configuration with multiple seeds and computes
several controls per cell. Three carry reference communities: email-Eu-core has
department labels, and com-DBLP and com-Amazon carry the overlapping communities
distributed with them. The mechanism measurements, which require only one
partition per graph and no restarts, run on a larger suite of seventeen networks
extending to 25 million edges.
All graphs are reduced to their largest connected component and treated as
simple and undirected.

\begin{table}[t]
\centering
\caption{The seven networks carrying the quality measurements, ordered by average
degree. Sizes are of the largest connected component. GT marks the three with
reference communities.}
\label{tab:networks}
\small
\setlength{\tabcolsep}{5pt}
\begin{tabular}{@{}lrrrc@{}}
\toprule
Network & $n$ & $m$ & $d_{\mathrm{avg}}$ & GT\\
\midrule
com-Amazon    & 334{,}863 &   925{,}872 &  5.53 & $\bullet$\\
ca-HepTh      &   8{,}638 &    24{,}806 &  5.74 & \\
com-DBLP      & 317{,}080 & 1{,}049{,}866 &  6.62 & $\bullet$\\
ca-CondMat    &  21{,}363 &    91{,}286 &  8.55 & \\
email-Enron   &  33{,}696 &   180{,}811 & 10.73 & \\
wiki-Vote     &   7{,}066 &   100{,}736 & 28.51 & \\
email-Eu-core &      986 &    16{,}064 & 32.58 & $\bullet$\\
\bottomrule
\end{tabular}
\end{table}

Real networks cannot answer questions about where a transition lies, because
their density cannot be varied while their community structure is held fixed. Two
families of synthetic benchmark do that. LFR
graphs~\cite{lancichinetti2008benchmark} with $10^4$ vertices and degree exponent
$2.1$ are generated at nominal average degrees of $10$, $25$, $50$, $100$ and
$200$, three graphs per configuration, matching the generator settings of the
original synthetic sweep. Stochastic block models and degree-corrected block
models are generated at the same degrees, with degree and mixing calibrated per
cell to the realized values of the LFR graphs, in two variants each so that
degree heterogeneity and community-size heterogeneity form a $2 \times 2$
factorial. Realized degree, realized mixing and degree coefficient of variation
are recorded per graph and all results are reported against realized values.

\paragraph{Sparsifiers}
Seven methods are tested, chosen so that a negative result cannot be attributed
to the selection. Two are the methods whose quality claims motivate the pipeline:
local sparsification, L-Spar~\cite{satuluri2011local}, which lets each vertex keep
the $\lceil d_v^{\,e} \rceil$ incident edges of highest neighbourhood Jaccard
similarity, and the degree-based sampler DSpar~\cite{liu2023dspar}. Three more
are the methods that the largest independent benchmark of sparsification ranks
highest at preserving clustering structure, namely K-Neighbor, Local Degree and
Local Similarity~\cite{chen2024demystifying, hamann2016structure}, so that the
arms include the strongest candidates chosen by someone other than us. L-Spar and
Local Similarity are distinct methods despite the similar names, and L-Spar
always refers to the method of Satuluri et al. The remaining two are a
connectivity-preserving backbone, a spanning structure topped up to the retention
target either at random or by Jaccard similarity, which cannot disconnect a graph
by construction and is included as an instrument rather than as a competitor.
Uniform random deletion supplies the no-information baseline wherever a selection
rule has to be shown to be doing work.

Two classes are excluded and the exclusions bound what follows. Node-removing
methods such as $k$-core peeling are not tested, because they score their
partition on a smaller and denser vertex set, which is a different comparison
requiring the coverage control of Section~\ref{sec:protocol:retention}. Learned
sparsifiers~\cite{wu2020gsgan} are not tested, because matching them on realized
retention requires retraining at each operating point.

\paragraph{Detection algorithms}
Four algorithms choose their own number of communities:
Leiden~\cite{traag2019louvain} and Louvain~\cite{blondel2008fast}, which maximize
\eqref{eq:modularity}, Infomap~\cite{rosvall2008maps}, which minimizes a
description length, and label propagation~\cite{raghavan2007near}. One takes the
number as an input: Metis~\cite{karypis1998metis}, through the \texttt{pymetis}
bindings, which also imposes a balance constraint on part sizes. Where Metis is
used, $k$ is fixed to the same value in the sparsified and unsparsified arms, so
the two partitions have identical community counts by construction and no
granularity effect is possible.

On real networks $k$ has to be chosen, and we report both defensible choices: the
community count a modularity optimizer returns on the original graph, and, on the
three networks with reference communities, the number of those communities. The
two differ by as much as a factor of sixteen and are reported separately
throughout.

Of the algorithms behind the founding accuracy claim, Metis is tested here on
every network. For Graclus and MLR-MCL we could find no maintained implementation
that we were able to run, so MLR-MCL, the algorithm that claim is largest about,
remains untested by us. This is the most consequential gap in the paper and
Section~\ref{sec:discussion} returns to it.

\paragraph{Operating points}
Sparsifiers are compared at matched realized retention, with targets of $0.5$ and
$0.2$. Four of the seven methods cannot reach the lower target on a sparse graph,
because each vertex retains at least one incident edge, so a third operating
point is added per network at the largest floor among the arms, ensuring one
aggressive point at which all seven are comparable. Table~\ref{tab:floors} gives
the floors, and they are a result in their own right rather than a nuisance
parameter, since they bound how far the pipeline can be pushed on the graphs
where it would need to pay off most.

\begin{table}[t]
\centering
\caption{Hard retention floors, as a fraction of edges. Below these values the
method cannot run at all, whatever target it is given. Networks are ordered by
average degree, which governs the floor. DSpar and uniform deletion have no
floor.}
\label{tab:floors}
\small
\setlength{\tabcolsep}{4pt}
\begin{tabular}{@{}lrrrrr@{}}
\toprule
Network & $d_{\mathrm{avg}}$ & Backbone & L-Spar & Loc.\ Deg. & K-Neigh.\\
\midrule
com-Amazon    &  5.53 & 0.362 & 0.301 & 0.352 & 0.330\\
ca-HepTh      &  5.74 & 0.348 & 0.276 & 0.347 & 0.317\\
com-DBLP      &  6.62 & 0.302 & 0.244 & 0.301 & 0.277\\
ca-CondMat    &  8.55 & 0.234 & 0.186 & 0.234 & 0.218\\
email-Enron   & 10.73 & 0.186 & 0.159 & 0.186 & 0.179\\
wiki-Vote     & 28.51 & 0.070 & 0.067 & 0.070 & 0.069\\
email-Eu-core & 32.58 & 0.061 & 0.053 & 0.061 & 0.060\\
\bottomrule
\end{tabular}
\end{table}

\paragraph{Replication}
Stochastic sparsifiers are run with independent sparsification seeds and
stochastic detectors with independent detection seeds; deterministic methods are
run once and reused. Baselines receive the two comparisons of
Section~\ref{sec:protocol:compute}, the compute-matched restarts and the best of
a fixed number. Where Metis is used the comparison is the worst sparsified run
against the best baseline run over ten partitioner seeds, which is the strictest
form available and the one used throughout for that detector.

\paragraph{Pre-registration}
Each experiment has a design document written before its code, fixing what is
measured, what is predicted, and what outcome would end that line of work.
Several predictions were ours and failed, and where a result below contradicts a
prediction we made, it is reported as such.

% ============================================================
% SECTION 4.2: WHAT THE ACCOUNTING COSTS. Assembled 2026-07-26 (late).
%
% RELOCATION, not rewriting. Every table and every number here is moved verbatim
% from sections/_mined/03_protocol_v1_full.tex, which held them while the
% protocol section was still carrying its own evidence. Only the connective
% prose is new, and its job is to make one argument out of what used to be six
% separate control demonstrations.
%
% Provenance of each figure is unchanged and was verified against CSVs earlier:
%   uniform table      exp_G/results.csv (arm=artifact1, 3 seeds, means)
%   63-cell table      exp_P/results.csv (dQ_naive, dQ_honest_vs_mean)
%   granularity table  exp_P/recovery.csv (com-DBLP, louvain)
%   compute inversion  exp_P/results.csv + bestof.csv; exp_C/results_summary.csv
%   permuted weights   exp_V/results.csv + null_arm.csv
%   retention audit    exp_J_figures/clip_fractions.csv; exp_C; exp_B
% STYLE: no em-dashes, no AI-characteristic phrasing (see STYLE.md).
% ============================================================

\subsection{What the accounting costs}\label{sec:experiments:accounting}

Each control in Section~\ref{sec:protocol} exists because a comparison can favour
the pipeline without any improvement having occurred. This subsection measures
how much each one is worth. The measurements are of the evaluation and not of the
sparsifiers, and they set the size of the correction that the rest of the paper
applies.

\paragraph{Scoring on the sparsified graph}
Table~\ref{tab:uniform} takes uniform random edge deletion, whose selection rule
uses no information whatever, and scores it both ways on three networks at two
retentions. Read off the sparsified graph it improves modularity in all six
cells. The same partitions scored on the original graph are worse than the
baseline in all six. The sign of the conclusion is opposite in every cell, and
the gap widens as retention falls, which is what
Section~\ref{sec:protocol:transfer} predicts: the sparser the graph, the higher
the modularity an optimizer can reach on it.

\begin{table}[t]
\centering
\caption{Uniform random edge deletion under the two accountings, Leiden, mean of
three seeds. $\Delta Q_{\mathrm{self}}$ is positive in every cell and
$\Delta Q_{\mathrm{transfer}}$ is negative in every cell. Uniform deletion uses
no structural information, so the inflation comes from the objective rather than
from any selection rule.}
\label{tab:uniform}
\small
\setlength{\tabcolsep}{4pt}
\begin{tabular}{@{}lccrr@{}}
\toprule
Network & $Q_{\mathrm{base}}$ & $\rho$ & $\Delta Q_{\mathrm{self}}$ &
$\Delta Q_{\mathrm{transfer}}$\\
\midrule
ca-CondMat  & 0.729 & 0.800 & $+0.0111$ & $-0.0164$\\
ca-CondMat  & 0.729 & 0.900 & $+0.0053$ & $-0.0064$\\
email-Enron & 0.611 & 0.800 & $+0.0052$ & $-0.0162$\\
email-Enron & 0.611 & 0.900 & $+0.0011$ & $-0.0086$\\
com-DBLP    & 0.826 & 0.800 & $+0.0078$ & $-0.0247$\\
com-DBLP    & 0.826 & 0.900 & $+0.0040$ & $-0.0098$\\
\bottomrule
\end{tabular}
\end{table}

The effect is not particular to that sparsifier or that detector.
Table~\ref{tab:accounting} covers 63 configurations spanning three detection
algorithms, seven networks and three sparsifier settings. Modularity read off the
sparsified graph is positive in 60 of them and reaches $+0.537$; the same
partitions scored on the original graph are positive in 15. Between the two
accountings the sign of the conclusion differs in 45 of the 63 configurations.
Eleven cells exceed a compute-matched baseline, all of them under label
propagation, and they are not a subset of the 15: three sit below the plain
baseline mean while beating a matched control that had drawn a single restart,
which is the failure the compute control addresses.

\begin{table}[t]
\centering
\caption{The two accountings across 63 configurations, seven networks by three
sparsifier settings under each algorithm. Both quantities are measured against
the mean of the plain restarts, the weakest of the baselines in
Section~\ref{sec:protocol:compute}; counts against the stricter best-of-$N$
baseline appear in Section~\ref{sec:experiments:negative}. A flip is a cell where
$\Delta Q_{\mathrm{self}} > 0$ and $\Delta Q_{\mathrm{transfer}} < 0$, so that
the two accountings disagree on the sign of the conclusion.}
\label{tab:accounting}
\small
\setlength{\tabcolsep}{4pt}
\begin{tabular}{@{}lcccc@{}}
\toprule
 & Infomap & Louvain & Lab.\ prop. & Total\\
\midrule
Cells                              & 21 & 21 & 21 & 63\\
$\Delta Q_{\mathrm{self}} > 0$     & 20 & 20 & 20 & 60\\
$\Delta Q_{\mathrm{transfer}} > 0$ & 6  & 0  & 9  & 15\\
Flips                              & 14 & 20 & 11 & 45\\
max $\Delta Q_{\mathrm{self}}$ & $+0.312$ & $+0.260$ & $+0.537$ & $+0.537$\\
\bottomrule
\end{tabular}
\end{table}

\paragraph{Granularity}
Table~\ref{tab:granularity} gives the anatomy of a recovery gain on com-DBLP
under Louvain. L-Spar at realized retention $0.503$ raises average best-match
$F_1$ from $0.102$ to $0.193$, an apparent gain of $+0.091$, while taking the
number of communities of at least three members from 220 to 10{,}947.
Partitioning the original graph to the same count reaches $0.300$. The finer
partition is worth more than twice what the sparsification appeared to be worth
and requires no sparsifier. The same pattern holds for the degree-based arm on
the same network, where the apparent change is negative and the matched control
is positive. Leiden reproduces the sparsified figure to four decimal places and
its own control, at $\gamma = 576$ and 10{,}812 communities, reaches $0.297$, so
the effect belongs to the modularity objective's default granularity and not to
one optimizer's search.

\begin{table}[t]
\centering
\caption{Granularity control on com-DBLP under Louvain, average best-match $F_1$
over reference communities of size at least three. $k_{\geq 3}$ counts
communities of that size; $\gamma$ is the resolution used on the \emph{original}
graph. Each sparsified arm is followed by its control, a partition of the
original graph at the same $k_{\geq 3}$. The floor is a size-matched random
partition at the arm's own $k_{\geq 3}$. Baseline and sparsified rows are means
over three to six seeds; controls are single runs.}
\label{tab:granularity}
\small
\setlength{\tabcolsep}{4pt}
\begin{tabular}{@{}lrrcc@{}}
\toprule
Condition & $k_{\geq 3}$ & $\gamma$ & $F_1$ & floor\\
\midrule
Original graph               & 220    & 1   & $0.102 \pm 0.005$ & 0.019\\
DSpar, $\rho = 0.500$        & 1{,}684  & 1   & $0.098 \pm 0.002$ & 0.049\\
\quad matched control        & 1{,}657  & 24  & $0.125$           & \\
L-Spar, $\rho = 0.503$       & 10{,}947 & 1   & $0.193 \pm 0.000$ & 0.090\\
\quad matched control        & 11{,}159 & 640 & $0.300$           & \\
\bottomrule
\end{tabular}
\end{table}

Where the counts cannot be brought together, we draw no recovery conclusion. The
comparability rule is $|\Delta k|/k < 25\%$, and the cells it excludes are
reported as excluded rather than compared.

\paragraph{Chance}
The floor column of Table~\ref{tab:granularity} carries the second artifact. A
random partition with the same community size distribution scores $0.019$ against
the reference communities at 220 communities and $0.090$ at 10{,}947, so a
comparison that reads $+0.091$ from a fragmenting sparsifier is reading a floor
that rose by $0.071$ underneath it. The same measurement on com-Amazon gives
$0.009$ at 231 communities and $0.057$ at 8{,}651, a $6.6\times$ rise from
fragmentation alone.

\paragraph{The compute budget}
The registered compute-matched control buys less than it appears to. In a sweep
of 48 cells over six networks, four retention targets and two samplers under
Leiden, it bought exactly two restarts in every cell; across the 63
configurations of Table~\ref{tab:accounting} it bought between 1 and 16, and 1 in
13 of the 18 cells on the two largest networks. Label propagation shows the
consequence. Under the compute-matched control it is the strongest of the three
algorithms in that table, with a mean change of $+0.009$; against the best of 20
plain restarts of the same baseline, 10 on the two largest networks, it is the
weakest, at $-0.106$. On email-Enron its two apparent gains of $+0.146$ and
$+0.161$ become $-0.068$ and $-0.053$. The compute-matched baseline drew four and
five restarts in those two cells and still never sampled the algorithm's good
mode, whose modularity is $0.523$ against the matched control's $0.308$.

\paragraph{Weighting credited to similarity}
On email-Enron, weighting each edge by $1 + J(u,v)$, with $J$ the Jaccard
similarity of the endpoint neighbourhoods, raises modularity measured on the
original graph to $0.6153 \pm 0.0034$ against a baseline of $0.6071$ and a
compute-matched best of $0.6051$, a gain of $+0.0102$ at $3.7$ times the pooled
standard deviation. Permuting those same weights across edges, which preserves
their distribution and destroys their placement, gives $0.6168 \pm 0.0020$, a
gain of $+0.0117$ at $6.0$ pooled standard deviations. The null reproduces the
gain and exceeds it. On a degree-preserving rewiring of ca-CondMat, where the
similarity signal is absent by construction, weighting still produces $+0.0023$
and clears the same significance test. What the weights contribute is
heterogeneity in the optimizer's search and not information about communities.

\paragraph{Nominal retention}
The parameter a sampler is given does not determine what it deletes. The scheme
in common use clips per-edge probabilities at one, and across six networks at
four targets between 12\% and 36\% of edges reach that clip, so the expected
number retained falls below the nominal $\alpha m$ and saturates: on
email-Eu-core the expected retention at $\alpha = 1.0$ is $0.665$, and on
wiki-Vote at $\alpha = 0.95$ it is $0.447$. Three samplers that describe
themselves by the same $\alpha$ realize different fractions. Sampling with
replacement gives $0.33$ to $0.52$ across seventeen networks at a nominal $0.8$,
the clipped no-replacement scheme gives $0.37$ to $0.68$ across those six
networks and four targets, and solving the scale factor by bisection on the same
cells gives $0.70$ to $0.95$. A comparison at equal nominal parameters is a
comparison at unequal deletion.

% ============================================================
% SECTION 4.3-4.5: THE NEGATIVE, ITS RIVAL EXPLANATIONS, AND THE EXCEPTION.
% Restructured 2026-07-26 (late) from sections/04_below.tex, whose full previous
% form is preserved at sections/_mined/04_below_v1_full.tex.
%
% Changes in this pass, all structural:
%   - promoted from a standalone section to subsections of Experiments
%   - the fixed-k result and the speed result become paragraphs of 4.3 rather
%     than subsections of their own, since each is one measurement
%   - the recovery subsection is deleted: it restated the com-DBLP anatomy that
%     4.2 now owns, so 4.3 points at Table 5 instead of retelling it
%   - "four ways the loss could be an artifact" keeps its own subsection because
%     its four legs need paragraph headings of their own
% No numbers changed in this pass.
% STYLE: no em-dashes, no AI-characteristic phrasing (see STYLE.md).
% ============================================================

\subsection{Where sparsification does not help}\label{sec:experiments:negative}

\paragraph{The objective}

Table~\ref{tab:negative} reports the result. Across 102 matched-retention cells,
seven sparsifiers on seven networks under a modularity optimizer, not one cell
improves modularity measured on the original graph, whether the baseline is the
compute-matched control or the best of five plain restarts. The best cell in the
entire grid loses $0.0105$.

\begin{table}[t]
\centering
\caption{Every matched-retention cell, Leiden, seven networks. $\Delta Q$ is
measured on the original graph, Eq.~\eqref{eq:transfer}, against the best of five
plain restarts. Counts are cells with a positive value. The arms differ in how
much they lose and not in whether they lose.}
\label{tab:negative}
\small
\setlength{\tabcolsep}{4pt}
\begin{tabular}{@{}lrrrr@{}}
\toprule
Arm & Cells & $\Delta Q > 0$ & Best cell & Worst cell\\
\midrule
K-Neighbor            & 18 & 0 & $-0.0105$ & $-0.667$\\
DSpar                 & 18 & 0 & $-0.0116$ & $-0.640$\\
L-Spar                &  9 & 0 & $-0.0117$ & $-0.407$\\
Backbone, random fill & 14 & 0 & $-0.0141$ & $-0.214$\\
Local Similarity      & 15 & 0 & $-0.0172$ & $-0.311$\\
Local Degree          & 14 & 0 & $-0.0288$ & $-0.278$\\
Backbone, Jaccard fill& 14 & 0 & $-0.0290$ & $-0.235$\\
\midrule
Total                 & 102 & \textbf{0} & $-0.0105$ & $-0.667$\\
\bottomrule
\end{tabular}
\end{table}

The sweep that isolates one method across its full retention range agrees. Local
L-Spar on the same seven networks loses between $0.014$ and
$0.404$ against a compute-matched baseline in every configuration, with
the loss growing as retention falls. Its three most aggressive configurations
cannot reach the retention they were given and run at the floors of
Table~\ref{tab:floors} instead.

\paragraph{Under a fixed-$k$ partitioner}
The results above concern
detectors that choose their own granularity. The condition under which
sparsification does help, established in this paper on synthetic benchmarks,
involves a detector that takes the number of communities as an input, so the
question is whether such a detector escapes the negative on these graphs.

It does not. Metis was run on all seven networks at the same operating points and
under the same protocol, with ten partitioner seeds per cell and the worst
sparsified run compared against the best baseline run. On the five networks with
average degree between $5.5$ and $10.7$, no cell is positive: 128 cells, seven
sparsifiers, and both ways of choosing $k$ on the two of those networks that
carry reference communities. The best cell loses $0.017$. Fixing
the community count removes the granularity artifact by construction and leaves
the loss exactly where it was.

The two conventions for choosing $k$ agree throughout, whether $k$ is taken from
the community count a modularity optimizer returns on the original graph or, on
the three networks with reference communities, from the number of those
communities. On com-Amazon the two values differ by a factor of sixteen, 306
against 5,000, and both give the same answer.

Above that range the picture changes, and Section~\ref{sec:experiments:boundary} takes it up:
on one of the two networks above it, four of the seven sparsifiers produce a gain
that survives the worst case over ten seeds. On the other, which is denser, none
does.

\paragraph{Recovery}
Agreement with reference communities behaves differently from the objective, and
the difference is instructive rather than encouraging. Raw agreement often rises,
and Table~\ref{tab:granularity} has already shown why: the sparsified partition is
finer, the measure rewards that, and the floor beneath it rises at the same time.
After both controls something still remains on two of the three labelled
networks, and Section~\ref{sec:experiments:exception} takes those up.

\paragraph{Speed}
The pipeline is not faster. Of the 63 configurations, four preserve quality in the
sense of coming within $0.005$ of the best-of-$N$ baseline, and the fastest of
those four runs at $0.66\times$ the speed of clustering the original graph
directly. Twelve configurations exceed $1.5\times$, and every one of them loses
modularity, by $0.015$ to $0.162$. Timing only the detection step and giving the
sparsifier away for free, the median speedup is $1.21\times$, $1.63\times$ and
$1.15\times$ by algorithm, well short of the factor of two that halving the edge
count suggests, because all three algorithms return more communities on the
sparsified graph and do more work per pass.

\subsection{Four ways the loss could be an artifact}\label{sec:experiments:artifact}

A negative result is only as good as the explanations it excludes. Four are
available, and each was tested rather than argued.

\paragraph{Fragmentation}
Sparsification breaks graphs apart, and the optimizer may be penalized for being
handed a shattered graph. The backbone arms test this directly, because they
retain a spanning structure and cannot increase the number of connected
components. They behave as designed: under the backbone with random fill no
detected community is a singleton in any of its fourteen cells, and the fraction
of vertices in communities of fewer than ten members never exceeds $1.8\%$,
where the shattering arms reach $86\%$ on the same networks. The conclusion does
not move. The backbone's best cell is $-0.0141$, worse than three of the arms that
fragment freely. Removing half the edges costs modularity on the original graph
whether or not the graph stays in one piece.

\paragraph{The choice of optimizer}
Everything above uses a modularity optimizer, so the objective and the algorithm
may be too closely coupled. Two sparsifiers across three further algorithms, one
flow-based, one a local majority rule and a second modularity optimizer, give the
63 configurations of Table~\ref{tab:accounting}. Against the best of $N$ plain
restarts, 60 are negative, with mean changes of $-0.082$ under Infomap, $-0.046$
under Louvain and $-0.106$ under label propagation. The three positive cells are
label propagation on one network whose unsparsified baseline is pathological: its
own best of twenty restarts reaches modularity $0.079$ where a modularity
optimizer on the same graph reaches $0.416$. Relieving a known failure mode of one
algorithm is not a benefit of sparsification.

\paragraph{A benign loss}
The loss might be harmless if sparsification sheds peripheral vertices while
leaving community cores intact. It does not. The fragments subdivide communities
without mixing them, but they cut the most embedded vertices as readily as the
least, and a resolution-matched partition of the original graph preserves cores
better in sixteen of twenty comparisons. Restricting the damage to leaf and
near-leaf vertices does not rescue it either.

\paragraph{The retention regime}
The loss might be an artifact of pushing the methods too hard. It is not: it is
already present at the mildest retention any of them can reach. On the three
sparsest networks that is about a third of the edges, by Table~\ref{tab:floors},
and every arm loses there. The regime in which sparsification would pay for itself
is not merely unprofitable on these graphs, it is unreachable.

One observation from the same arms does not belong to this list, because it
supports no rival explanation. The backbone's two fills differ only in which edges
top up the spanning structure, random or highest Jaccard similarity, and random
fill scores better in ten of the fourteen paired cells. In this regime the
similarity signal subtracts.

\subsection{The exception}\label{sec:experiments:exception}

Two networks carry a result that clears the matched control of
Section~\ref{sec:protocol:granularity} and the chance floor.

On com-Amazon, L-Spar at realized retention $0.542$
reaches average best-match $F_1$ of $0.402$ against a baseline of $0.142$. The
matched control does not remove it: partitioning the original graph to the
same community count gives $0.319$, and deliberately over-matching by 37 per cent
more communities still gives only $0.372$. The chance floor does not remove it
either, and subtracting the floor widens the margin rather than narrowing it,
because the floor rises with the community count. A second instance
appears in the same experiments from an unrelated method, Local Degree, on
com-Amazon and on com-DBLP. The two baselines separate on com-DBLP: that instance
clears the matched control, at $0.163$ against $0.123$ at the same community
count, and does not clear the sweep, whose best point on the untouched graph
reaches $0.397$.

What removes both is a different comparison. Both non-modularity algorithms reach
higher agreement on com-Amazon without any sparsification at all, $0.465$ under
Infomap and $0.480$ under label propagation, against the $0.402$ that the
sparsified modularity pipeline attains, and the same holds on com-DBLP at
$0.326$ and $0.382$ against $0.163$. The effect is real, and it is the
modularity objective's default granularity being partially repaired by shattering
the graph, in a metric that rewards fine partitions. It is not a capability that
sparsification confers.

Two things about this cell are worth carrying forward. It coexists with a
modularity loss of $0.079$ in the same configuration, so the objective and the
recovery measure move in opposite directions on the same partition of the same
graph. And it is the reason the negative result of this section is stated as it
is: no benefit on the objective anywhere, and no recovery benefit that survives
comparison against a better algorithm run on the untouched graph. Section
\ref{sec:discussion} takes up the dissociation as an open phenomenon.

% ============================================================
% SECTION 4.6: THE BOUNDARY. Written fresh 2026-07-26 (late) from the CSVs.
% This REPLACES sections/05_boundary.tex, which Mohammad retired: it predated
% the protocol section, and three of its numbers did not survive checking.
%
% Numbers verified in this pass, with what changed from the retired section:
%   LFR counts by degree      exp_AA via exp_AC/analyse.py LFR reference block
%   seed robustness           RECOMPUTED from exp_AA/metis_noise.csv. The retired
%                             text said the effect clears noise "by a factor of
%                             between two and sixty". True range at d>=49 is
%                             -1.32 to 5.79 in modularity and -1.88 to 14.36 in
%                             AMI, and several cells at retention 0.5 are
%                             negative. Restated per retention point.
%   "harmful on both measures below 25"  CONTRADICTED by exp_AA's own table
%                             (6/12 positive on AMI at 24.6). Removed.
%   speed                     RECOMPUTED from exp_AA/results_armB.csv: median
%                             pipeline speedup 0.75x at d=10 falling to 0.26x at
%                             d=200. The retired text's 1.49x / 1.01x / 0.58x
%                             does not match arm B.
%   edge-budget control       DSpar is positive in 2/12 and 3/12 cells at the two
%                             highest degrees even on LFR, and on the one real
%                             network above the transition it reproduces the gain
%                             outright (exp_AD). "Does not reproduce at any
%                             degree" is retired.
%   Leiden arm                exp_AA/results_armA.csv: 28/144 lspar cells beat
%                             best-of-5, none above +0.005; resolution-matched
%                             control wins 104/144.
%   generalization            exp_AC/SUMMARY.md and analysis_full.txt
%   real networks             exp_AD/results.csv
%   denoising                 exp_AA/results_armC.csv, leiden+lspar 4/4 monotone
% STYLE: no em-dashes, no AI-characteristic phrasing (see STYLE.md).
% ============================================================

\subsection{The boundary}\label{sec:experiments:boundary}

Everything so far is measured on one side of a line. The networks in
Table~\ref{tab:networks} have average degrees between $5.5$ and $32.6$, and the
detectors used for most of it choose their own granularity. Both conditions in
this paper's claim are therefore held at one setting, and a negative result
obtained that way says nothing about where the setting stops mattering. This
subsection varies both.

\paragraph{What the setting requires}
Three properties are needed. Density must vary over an order of magnitude while
community structure is held fixed, which real networks cannot provide. Community
membership must be known rather than inferred. And the granularity artifact of
Section~\ref{sec:experiments:accounting} must be removed by construction rather
than by control, which a detector taking $k$ as an input does: both arms return
the same number of communities by definition.

\paragraph{The transition}
Table~\ref{tab:boundary} reports, for LFR graphs at each average degree, the
number of configurations in which L-Spar improves modularity measured on the
original graph under Metis, and the number in which it improves chance-corrected
agreement with the planted labels.

\begin{table}[t]
\centering
\caption{Metis with $k$ pinned to the planted community count, LFR, realized
mixing near $0.6$. Counts are configurations out of twelve in which L-Spar
improves the quantity named, with the mean change beside them. Modularity is
measured on the original graph. Because $k$ is an input, both arms return
identical community counts and no granularity effect is possible.}
\label{tab:boundary}
\small
\setlength{\tabcolsep}{5pt}
\begin{tabular}{@{}lcccc@{}}
\toprule
$d_{\mathrm{avg}}$ & $\Delta Q > 0$ & mean $\Delta Q$ &
$\Delta\mathrm{AMI} > 0$ & mean $\Delta\mathrm{AMI}$\\
\midrule
11.2  & 0/12  & $-0.103$ & 0/12  & $-0.216$\\
24.6  & 0/12  & $-0.046$ & 6/12  & $-0.072$\\
49.4  & 8/12  & $-0.016$ & 9/12  & $+0.042$\\
101.9 & 10/12 & $+0.006$ & 12/12 & $+0.116$\\
220.9 & 12/12 & $+0.009$ & 12/12 & $+0.093$\\
\bottomrule
\end{tabular}
\end{table}

The two measures cross at different places. Recovery turns positive first, with
half the configurations improving at average degree $24.6$ and three quarters at
$49.4$, while the objective is still negative on average there and only becomes
positive in the mean at $101.9$. The location of the crossing was not chosen by
us and was not known to us when the experiment was designed: it coincides with
the degree at which the original authors reported their method beginning to
overtake the unsparsified baseline.

The size of the effect should be read with the counts. At the degree where the
objective crosses, the mean change is still negative and the positive cells are
worth thousandths of a modularity point. The recovery gains are an order of
magnitude larger, reaching a mean of $+0.116$.

\paragraph{Run-to-run variation}
Metis is not deterministic across option seeds, so the decisive cells were
repeated with ten seeds per graph and the worst sparsified run compared against
the best unsparsified run. The comparison separates the two measures again. On
recovery it is comfortable: at aggressive retention the worst case clears the
baseline seed spread by factors of $4.9$ to $19.8$, and it clears it at every
degree at or above $49$. On the objective it is not: the same cells reach factors
of only $7.8$ down to $-0.4$, so one of them fails to clear at all, and at the
mildest retention the cells at $201$ and $261$ are negative in the worst case on
both measures. The recovery result survives a hostile reading of the seeds; the
modularity result survives at the aggressive operating points with one exception.

\paragraph{Generality across generators}
LFR graphs are one family with one degree distribution, and the mechanism runs
through a degree-derived or similarity-derived score, so the family could be
doing the work. Repeating the arm on stochastic block models and on
degree-corrected block models, calibrated cell by cell to the LFR graphs'
realized degree and mixing, reproduces the sign flip on both. It also refutes the
explanation we expected. We predicted that a degree-homogeneous generator would
push the transition to higher density, since a degree-based selector has less to
work with, and that the degree-corrected variant would sit closer to LFR. Neither
holds. All four block-model variants cross at essentially the same place while
their degree coefficient of variation spans an order of magnitude: at the degree
where the transition sits it runs from $0.14$ to $1.47$ across the four
generators. Across sixty generator-by-degree cells the gain correlates with
average degree at Spearman $+0.65$ and with degree variability at $-0.26$. Degree
heterogeneity is not the operative variable.

\paragraph{Real networks}
The same arm on the seven real networks places the transition without resolving
what sits on the far side of it. Below average degree $11$ there is nothing: 128
cells across five networks, seven sparsifiers, and both ways of choosing $k$ on
the two of those networks that carry reference communities, with not
one positive in the worst case over ten seeds. On wiki-Vote, at average degree
$28.5$, four of the seven sparsifiers produce a worst-case positive gain, the
largest being $+0.0124$. On email-Eu-core, at $32.6$, none does: two arms are
positive in the mean and every worst case is negative. The transition on real
graphs therefore lies between $10.7$ and $28.5$, and above it the effect is
present on one of the two networks we can test.

\paragraph{What produces the gain}
On the benchmark graphs the selection rule matters: degree-based and uniform
sparsification at identical retention are negative at every degree except the two
highest, where DSpar turns positive in 2 and 3 of twelve cells against L-Spar's
10 and 12. On wiki-Vote that separation disappears. The four arms that score
individual edges locally all produce the gain, DSpar among them at $+0.0076$
against L-Spar's $+0.0124$, and the two connectivity-preserving backbones destroy
it. What the gain requires is a rule that scores edges one at a time; which
signal the rule uses matters less than we expected, and we do not identify the
property that separates Local Degree, which is also local and is the worst arm on
that network.

\paragraph{Why density}
The denoising account makes a testable prediction. If sparsification helps by
deleting edges that do not belong to the generative structure, its benefit should
grow with the proportion of such edges. Injecting uniformly random edges into
benchmark graphs while holding the planted labels fixed is consistent with it for
similarity-based selection: recovery gain rises with the injected
fraction in three of the four configurations tested and rises with one reversal
in the fourth, and at average degree
$50$ with $20\%$ retention it grows from $+0.061$ with no injected noise to
$+0.185$ when the number of spurious edges equals the number of real ones. The
gain does not vanish at zero noise, so denoising adds to an effect already
present rather than being its sole cause, and a benchmark's noise is random by
construction whereas a measurement process need not be. The experiment
establishes the direction, not the magnitude in the field.

\paragraph{The same graphs, the other detector family}
Running a free-granularity optimizer on the identical graphs gives a different
answer, which is what makes the boundary a boundary rather than a threshold in
density alone. Honest modularity gains appear in 28 of 144 configurations and
none exceeds $+0.005$, within the spread of five random restarts. Recovery gains
do appear and grow with degree, but they do not survive the granularity control:
partitioning the original graph at a resolution tuned to the sparsified
partition's community count recovers the planted communities as well or better in
104 of the 144 configurations where that comparison is defined. What looked like
a benefit of sparsification is the benefit of a finer partition, available
without touching the graph.

\paragraph{Three things the positive result does not say}
Sparsification repairs a weaker detector rather than producing the best available
partition. At average degree $49$, Metis with L-Spar reaches AMI $0.75$ against
its own baseline of $0.62$, a real improvement to that algorithm, while a
resolution-tuned optimizer on the untouched graph reaches $0.94$ on the same
graphs. The defensible statement is that similarity-based selection repairs a
fixed-$k$ cut partitioner on dense graphs, not that it improves community
detection.

The quality gain is not a speed gain in our implementation. Computing exact
pairwise similarity costs more than it saves, and the deficit grows with density
because the similarity computation grows faster than the detection it
accelerates: the median end-to-end speedup falls from $0.89\times$ at average
degree $11$ to $0.18\times$ at $221$. The original speedups were obtained with an
approximate similarity computation reported as two orders of magnitude cheaper,
measured against clustering runs lasting hours, so both cost bases differ from
ours by enough to explain the discrepancy without disagreement about what should
be counted.

One part of the original report does not reproduce. At high mixing, where the
planted structure is barely detectable and baseline recovery sits between $0.03$
and $0.17$, no gain appears at any degree: across 104 configurations L-Spar
improves modularity in 9, and the mean change is negative at every degree. The
original sweep reports its largest improvement in that regime. Whether this is a
genuine non-reproduction or a floor effect of the metrics we use cannot be
settled at that mixing level, and we report it unresolved.

% ============================================================
% SECTION 4.7: MECHANISM. Written 2026-07-26 (late) after the audit.
%
% SCOPE, established by the audit and easy to overstate: this explains the
% DEGREE-BASED channel, not the boundary. The boundary's mechanism is the
% denoising evidence in 4.6. PAPER_STRUCTURE.md called the old section
% "Mechanism: why the boundary is there", which was wrong.
%
% Numbers, all recomputed this session:
%   exp_B/results.csv        17 networks; dQ_fixed positive 17/17 real and
%                            17/17 null, null larger in 15; delta larger on the
%                            null in 16/17, median ratio 1.599
%   exp_M/probe_results*.csv decomposition on the NULL rows (the terms are
%                            defined relative to the real arm, so real rows are
%                            zero by construction): sorting positive 16/16,
%                            median +0.5743, range +0.08 to +3.42; identity
%                            term_sum = log_ratio to 1.7e-14; facebook-combined
%                            excluded, its real delta is -0.0014
%   exp_Q/verdicts.csv       pfix arm only: delta moves 37 to 1322 neutral sd,
%                            sign change 5/5; re-detection AMI 0.824-0.944 on the
%                            UP ends and 0.410-0.540 on the DOWN ends; the
%                            literal arm is excluded because it destroys the
%                            partition (AMI <= 0.038, p_intra -> 0)
% STYLE: no em-dashes, no AI-characteristic phrasing (see STYLE.md).
% ============================================================

\subsection{Mechanism of the degree-based channel}\label{sec:experiments:mechanism}

A degree-based sparsifier keeps an edge with probability proportional to a score
computed from its endpoints' degrees. Whether that helps or harms a partition
depends on one quantity: the difference between the mean score on edges inside
communities and the mean score on edges between them. Write it $\delta$. When
$\delta > 0$ the sampler removes between-community edges preferentially, which
raises the intra-community fraction of any fixed partition and therefore its
modularity. This subsection establishes what controls $\delta$, and it concerns
the degree-based channel only. The mechanism behind the boundary of
Section~\ref{sec:experiments:boundary} is the denoising evidence given there.

\paragraph{$\delta$ is not evidence about communities}
Two quantities in this literature are read as showing that a sparsifier clarifies
community structure: $\delta$ itself, and the modularity change of a fixed
partition when the graph is sparsified. Both were measured on seventeen real
networks ranging from $10^3$ to $3.8 \times 10^6$ vertices, and on
degree-preserving rewirings of each, built by double edge swaps, which hold every
vertex degree exactly and destroy the communities.

Neither survives. The fixed-partition modularity change is positive on all
seventeen real graphs; on their rewirings it is positive on all seventeen as
well, and larger than the real value on fifteen. The separation $\delta$ is
larger on the rewiring than on the real graph in sixteen of seventeen cases, with
a median ratio of $1.60$. Real community structure suppresses the separation
rather than producing it. Neither quantity is used as evidence anywhere else in
this paper, and the fixed-partition measure in particular is the same
self-scoring comparison that Section~\ref{sec:protocol:transfer} rejects.

\paragraph{Where the null's advantage comes from}
The ratio between the null's separation and the real graph's decomposes exactly.
Writing $\delta = r_{pb}\, \sigma_s / \sqrt{p(1-p)}$, with $r_{pb}$ the
point-biserial correlation between an edge's score and whether it lies inside a
community, $\sigma_s$ the spread of the scores and $p$ the intra-community edge
fraction, gives
\begin{equation}\label{eq:delta-decomp}
\log \frac{\delta_{\mathrm{null}}}{\delta_{\mathrm{real}}}
= \underbrace{\log \frac{r_{pb}^{\mathrm{null}}}{r_{pb}^{\mathrm{real}}}}_{\text{sorting}}
+ \underbrace{\log \frac{\sigma_s^{\mathrm{null}}}{\sigma_s^{\mathrm{real}}}}_{\text{scale}}
+ \underbrace{\tfrac{1}{2}\log \frac{p^{\mathrm{real}}(1-p^{\mathrm{real}})}
{p^{\mathrm{null}}(1-p^{\mathrm{null}})}}_{\text{granularity}} ,
\end{equation}
an identity that holds in our measurements to $1.7 \times 10^{-14}$. The sorting
term accounts for the difference. It is positive on all sixteen networks where
the ratio is defined, with a median of $+0.57$ and a range from $+0.08$ to
$+3.42$, while the scale and granularity terms are small and mostly negative. The
seventeenth network is excluded because its real separation is $-0.0014$, which
makes the ratio undefined. What the rewiring changes is not the spread of the
scores or the number of communities but how well the score sorts within-community
edges from between-community ones, and it sorts them better once the communities
are gone.

\paragraph{What controls $\delta$}
The correlational evidence points at where high-degree edges sit relative to
community boundaries, but a graph statistic that co-varies with $\delta$ across
networks does not establish that moving it moves $\delta$. We therefore steer it
directly. Starting from a real graph and its detected partition, edges are
rewired to increase or decrease the concentration of high-degree-product edges on
community boundaries, while the degree sequence, the partition, the
intra-community edge fraction and the modularity are all held exactly fixed. Only
the placement of hub edges relative to the boundaries changes.

$\delta$ follows. Across five networks it moves between 37 and 1322 standard
deviations of the unsteered chains, in the direction it is steered, and it
changes sign in all five when steered downward. The largest excursion takes
ca-GrQc from $+0.087$ to $-0.875$ with its degree sequence, partition,
intra-edge fraction and modularity unchanged.

Two limits on that result are worth stating. The steered graphs remain
recognizable in the upward direction, where re-detecting communities from scratch
recovers the original partition at AMI $0.82$ to $0.94$, and less so downward,
where re-detection agrees at $0.41$ to $0.54$. And a looser variant of the same
procedure, which does not hold the intra-community fraction fixed, drives that
fraction to zero and destroys the partition altogether, with re-detection AMI
below $0.04$; nothing about the mechanism can be concluded from that variant and
we exclude it.

\paragraph{An explanation that fails}
Triangle structure is the natural alternative: edges inside communities close more
triangles, so a score correlated with triangle participation would produce
$\delta > 0$ without any reference to degree placement. The steering experiment
rules it out. Intra-community triangle density falls in both directions of the
steering, from $11.7$ to $9.0$ upward and to $11.4$ downward on ca-GrQc and
similarly elsewhere, while $\delta$ moves in opposite directions in the two arms.
A quantity that decreases while the target rises in one arm and falls in the
other does not control the target.

\paragraph{What this explains}
The degree-based channel is mechanical. It is governed by where high-degree-product
edges sit relative to community boundaries, a property that a degree-preserving
rewiring can reproduce and often exaggerate, and it is therefore not a route by
which a degree-based sparsifier clarifies community structure. That is consistent
with the negative territory of Section~\ref{sec:experiments:negative}, where the
degree-based arm behaves like the others, and with the finding in
Section~\ref{sec:experiments:boundary} that above the transition the gain does not
depend on which local scoring rule is used.


% \section{The Boundary}\label{sec:boundary}
% % ============================================================
% SECTION 5: THE BOUNDARY. Written clean 2026-07-26 from exp_AA.
% Evidence: PAPER_RESTRUCTURE/exp_AA_satuluri_regime/{SUMMARY.md,
%   results_armA.csv, results_armB.csv, results_armC.csv,
%   metis_noise.csv, lfr_generation.csv}
% Weight: HEAVY, equal to the whole negative-territory section
% (PAPER_STRUCTURE.md). This carries the paper's novel finding.
% STYLE: no em-dashes, no AI-characteristic phrasing (see STYLE.md).
% ============================================================

Everything to this point has been measured on real networks whose average degree
lies between 5.5 and 32.6, and on that territory the answer is uniformly
negative. That range is not a neutral choice. It is the range in which
sparsification for community detection is normally studied, and it lies entirely
below the one quantitative boundary the founding work
names~\cite{satuluri2011local}: a synthetic sweep in which local similarity
sparsification ``actually outperforms the original clustering starting from
degree 50.'' A negative result confined to one side of a boundary answers half
the question. This section supplies the other half.

\subsection{Design}

Three properties are required of the setting. Average degree must vary over an
order of magnitude while everything else is held fixed. Community labels must be
known rather than inferred. And the granularity artifact that accounts for most
of the apparent gains in this literature must be removable by construction rather
than by control.

Synthetic benchmarks with planted communities supply the first two. A fixed-$k$
partitioner supplies the third. If the number of communities is an input to the
algorithm, then the sparsified and unsparsified arms return the same number of
communities by definition, and no difference between them can be attributed to
one partition being finer than the other.

We generate LFR benchmark graphs~\cite{lancichinetti2008benchmark} with
$n = 10^4$ nodes and degree exponent $2.1$, matching the generator settings of
the original sweep, at nominal average degrees of $10$, $25$, $50$, $100$ and
$200$, at two mixing levels, with three graphs per configuration. Realized degree
and mixing are recorded per graph, and all results are reported against realized
values. Three sparsifiers are applied at matched realized retention: local
similarity sparsification, degree-based sparsification, and uniform random edge
deletion. Detection is performed both by a free-granularity modularity optimizer
and, through the \texttt{pymetis} bindings, by Metis, a fixed-$k$
balance-constrained cut partitioner from the family the original accuracy claim
concerns. Quality is scored on the original graph throughout. Recovery is
measured against the planted labels with chance-corrected indices and a
size-matched random-partition floor.

\subsection{The threshold}

Table~\ref{tab:boundary} reports, for each average degree, the number of
configurations in which similarity-based sparsification improves modularity
measured on the original graph, and the number in which it improves
chance-corrected agreement with the planted communities, out of twelve
configurations per degree.

\begin{table}[t]
\centering
\caption{Fixed-$k$ partitioner (Metis), realized mixing $\approx 0.6$. Counts are
configurations out of twelve in which similarity sparsification improves the
quantity named, with the mean change beside them. Modularity is measured on the
\emph{original} graph. Because $k$ is an input to the partitioner, the sparsified
and unsparsified arms return identical community counts, so no granularity effect
is possible.}
\label{tab:boundary}
\begin{tabular}{@{}lcccc@{}}
\toprule
Avg.\ degree & $\Delta Q > 0$ & mean $\Delta Q$ & $\Delta\mathrm{AMI} > 0$ & mean $\Delta\mathrm{AMI}$\\
\midrule
11.2  & 0/12  & $-0.103$ & 0/12  & $-0.216$\\
24.6  & 0/12  & $-0.046$ & 6/12  & $-0.072$\\
49.4  & 8/12  & $-0.016$ & 9/12  & $+0.042$\\
101.9 & 10/12 & $+0.006$ & 12/12 & $+0.116$\\
220.9 & 12/12 & $+0.009$ & 12/12 & $+0.093$\\
\bottomrule
\end{tabular}
\end{table}

Below average degree $25$, the method is harmful on both measures in every
configuration tested. At $49.4$ it becomes beneficial in two thirds of them. By
$101.9$ it is beneficial in all twelve configurations on recovery and ten of
twelve on the objective. The strongest individual cells occur at aggressive
retention: at average degree $49.4$ and $15\%$ retention, the partition obtained
from the sparsified graph scores $+0.005$ modularity and $+0.152$ AMI above the
partition obtained from the original graph, and at $220.9$ the modularity gain
reaches $+0.012$.

The location of this transition was not chosen by us, and it was not known to us
when the experiment was designed. It coincides with the degree at which the
original authors reported their method beginning to overtake the unsparsified
baseline.

\subsection{Three ways the result could be spurious}

\paragraph{Run-to-run variation}
Metis is not deterministic across option seeds. We repeated the decisive cells
with ten partitioner seeds per graph and compared the \emph{worst} sparsified run
against the \emph{best} unsparsified run. The gain survives that comparison at
every degree at or above $49.4$, with worst-case modularity differences of
$+0.001$, $+0.005$ and $+0.008$, and worst-case AMI differences of $+0.134$,
$+0.102$ and $+0.064$. Seed standard deviations are $0.002$ in modularity and
$0.011$ in AMI, so the effect exceeds the noise it must clear by a factor of
between two and sixty.

\paragraph{The edge budget rather than the selection}
Any sparsifier at $15\%$ retention produces a smaller graph that is denser within
communities, and a fixed-$k$ partitioner may simply behave better on such a
graph. Degree-based and uniform random sparsification, applied to the same graphs
at the same realized retention, do not reproduce the effect at any degree. Their
modularity changes run from $-0.145$ to $-0.003$, and their recovery changes from
$-0.480$ to $+0.032$. What produces the gain is the similarity signal used to
choose which edges to keep, not the number of edges kept.

\paragraph{Granularity}
The partitioner takes the community count as an input, and both arms are run with
the same value, so the two partitions being compared have identical community
counts by construction. This is the artifact that no amount of matching fully
excludes on a free-granularity detector, and here it is excluded structurally. We
also record cluster-size balance and find it slightly improved rather than
degraded under sparsification, so the gain is not purchased by producing more
uneven communities.

\subsection{The same graphs, a different algorithm family}

Running a free-granularity modularity optimizer on the identical graphs gives a
different answer. Modularity gains never exceed the spread of a handful of random
restarts. Recovery gains do appear, and they grow with degree, reaching $+0.096$
AMI, but they do not survive the granularity control. Partitioning the original
graph at a resolution tuned to match the sparsified partition's community count
recovers the planted communities as well or better in twenty-one of the
twenty-nine configurations where the comparison is defined. What appeared to be a
benefit of sparsification is the benefit of a finer partition, obtainable without
touching the graph.

The variable that decides the outcome is therefore the interaction of density
with the detection algorithm's degrees of freedom, not the sparsifier and not the
graph alone. A partitioner constrained to return a fixed number of balanced
groups is helped, above a density threshold, by having its misleading edges
removed. An optimizer free to choose its own granularity already has a cheaper
route to the same improvement, and takes it.

\subsection{Why density matters: the denoising account}

The original authors attributed their gains to the removal of spurious edges,
noting that their strongest results came from an assay network known to contain
many false-positive interactions. That explanation makes a testable prediction.
If sparsification helps by deleting edges that do not belong to the generative
structure, then its benefit should grow with the proportion of such edges.

We inject uniformly random edges into benchmark graphs while holding the planted
labels fixed, and measure recovery against those labels. The prediction holds.
Recovery gain rises monotonically with the injected fraction in every tested
configuration for similarity-based sparsification, with rank correlations of
unity between noise level and gain. At average degree $50$ and $20\%$ retention,
the gain grows from $+0.062$ with no injected noise to $+0.184$ when the number
of spurious edges equals the number of real ones. Degree-based sparsification
shows no such response, which is consistent with its selection signal being a
function of degree rather than of local structure.

Two qualifications apply. The gain does not vanish when no noise is injected, so
denoising adds to an effect that is already present rather than being its sole
cause. And a benchmark graph's noise is random by construction, whereas the
spurious edges produced by a real measurement process need not be. The experiment
establishes the direction of the mechanism, not its magnitude in the field.

\subsection{What the positive result does not say}

Two facts constrain how this finding should be read, and both come from the same
experiments.

Sparsification repairs a weaker detector rather than producing the best available
partition. At every degree tested, a resolution-tuned modularity optimizer
running on the untouched graph recovers the planted communities better than any
sparsified pipeline. At average degree $49.4$ it reaches $0.956$ AMI, against
$0.754$ for the sparsified fixed-$k$ pipeline. Sparsification lifts the fixed-$k$
partitioner from $0.596$ to $0.754$, a real and substantial improvement to that
algorithm, and still leaves it behind an algorithm that never needed
sparsification. The defensible statement is that similarity sparsification
repairs a fixed-$k$ cut partitioner on dense graphs, not that it improves
community detection.

In our implementation the quality gain is also not a speed gain. Computing exact
pairwise similarity costs more than it saves. End to end, the pipeline runs at
$1.49\times$ the baseline at average degree $10$, $1.01\times$ at $100$, and
$0.58\times$ at $200$, which is slower than not sparsifying at all, even though
detection alone accelerates by up to $7.2\times$. The original speedups were
obtained with an approximate similarity computation reported as two orders of
magnitude cheaper, measured against clustering runs that took hours. Both cost
bases differ from ours by enough to account for the discrepancy without any
disagreement about what should be counted.

One part of the original report does not reproduce. At high mixing, where the
planted structure is nearly absent and baseline recovery falls to between $0.03$
and $0.19$, we find no gain under either detector at any degree, whereas the
original sweep reports its largest improvement in that regime. Whether this is a
genuine non-reproduction or a floor effect of the metrics we use cannot be
settled at that mixing level, and we report it unresolved.

%



\section{Discussion}\label{sec:discussion}
% TODO : the dissociation; what the boundary means for practice.

%
% CONCLUSION : WRITTEN LAST.
%
\section{Conclusion}\label{sec:conclusion}
% TODO (last).

\bibliographystyle{IEEEtran}
\bibliography{refs}

\end{document}
